\documentclass[12pt,a4paper]{article}

\usepackage{cmap}
\usepackage[T2A]{fontenc}
\usepackage[utf8]{inputenc}
\usepackage[english,russian]{babel}
\usepackage[a4paper,left=2.5cm,right=2cm,top=2cm,bottom=2cm]{geometry}
\usepackage{booktabs}
\usepackage{array}
\usepackage{enumitem}
\usepackage[table]{xcolor}
\usepackage{hyperref}

\definecolor{neuro}{RGB}{108,92,231}
\definecolor{inknavy}{RGB}{35,41,70}

\hypersetup{
  colorlinks=true,
  linkcolor=neuro,
  urlcolor=neuro,
  pdftitle={Токены и нейросети — правила игры},
  pdfauthor={Токены и нейросети}
}

\setlength{\parindent}{0pt}
\setlength{\parskip}{6pt}

\newcommand{\term}[1]{«\textbf{#1}»}

\title{\textbf{\textcolor{inknavy}{Токены и нейросети}}\\[4pt]
\large правила игры-соревнования}
\date{}

\begin{document}
\maketitle
\vspace{-2.5em}

\section*{Легенда}

Идёт 2052 год. Деньги, нефть и недвижимость окончательно вышли из моды:
единственная твёрдая валюта человечества --- \emph{токен}, сгенерированный
большой языковой моделью. Центробанки измеряют инфляцию в токенах в секунду,
пенсию индексируют по курсу эмбеддинга, а на чёрном рынке из-под полы
торгуют несжатыми логитами.

Вы --- совет директоров ИИ-корпорации. Ваш кластер круглосуточно генерирует
токены, но то же самое делают и конкуренты, у которых, по слухам,
«почти AGI, осталось дообучить». Единственный путь к господству ---
\emph{архитектурные прорывы}: их совершают учёные-исследователи,
они же --- задачи с сайта informatics. Внедрённый прорыв приводит
в экстаз венчурных инвесторов (\emph{«Это же трансформер с душой!
Берём!»}), и они засыпают компанию деньгами и видеокартами.

Правда, пока R\&D-отдел возится с прорывом, он нагло откусывает GPU
у продакшена --- генерация замедляется. А за отдельную плату штатный
инсайдер готов вынести с чужого кластера \emph{утечку весов} (тест
к задаче). Инсайдер берёт предоплату, работает без гарантий и NDA
подписывает не глядя.

К дедлайну, объявленному Всемирной Комиссией по Выравниванию, побеждает
корпорация с самым большим счётом токенов. Остальных ждёт слияние,
поглощение и переобучение в чат-бота службы поддержки.

\section*{Как всё устроено}

\begin{enumerate}[leftmargin=1.6em,itemsep=2pt]

\item \textbf{Стартовый капитал.} Каждая корпорация начинает с одинаковым
счётом \term{Токенов} и одинаковой \term{Мощностью} кластера. Кластер
не останавливается ни на секунду: каждую секунду игрового времени счёт
пополняется на величину текущей \term{Мощности}. Даже если вся команда
ушла на стендап.

\item \textbf{Каталог прорывов.} Всем корпорациям доступен один и тот же
набор задач: несколько уровней сложности, в каждом уровне --- несколько
задач. Но порядок задач внутри уровня у каждой корпорации свой, секретный
и перемешанный: подглядывать в дашборд конкурента бессмысленно ---
их ячейка № 5 и ваша ячейка № 5 --- почти наверняка разные задачи.
Промышленный шпионаж через плечо не работает, мы проверяли на проде.

\item \textbf{Покупка прорыва (задачи).} Постановка задачи засекречена,
пока исследование не оплачено: платите \term{Цену задачи} со счёта ---
получаете ссылку на условие. Купить задачу команда может \emph{сама},
кнопкой на своей странице. Грант возврату не подлежит: потраченные
на покупку токены не возвращаются \emph{никогда}, даже при наличии
презентации на сорок слайдов.

\item \textbf{Плата за прогресс.} С момента покупки и до сдачи прорыва
R\&D-отдел откусывает мощности: \term{Мощность} снижается
на \term{Нагрузку от задачи}. Хотите купить всё и сразу --- считайте:
если токенов на счету меньше цены, а кластер и так дымится,
финансовый отдел сделку не одобрит.

\item \textbf{Утечка весов (тесты).} Любое количество,
по \term{Цене подсказки} за штуку, приобретается через ведущего ---
инсайдер работает только с директором лично. Токены за утечку
не возвращаются, качество весов обсуждению не подлежит (пункт 1
джентльменского соглашения с инсайдером).

\item \textbf{Внедрение прорыва (решение задачи).} Решение сдаётся
на informatics. Как только вердикт --- \emph{OK} или \emph{Зачтено},
корпорация автоматически получает всё и сразу:
\begin{itemize}[itemsep=0pt]
	\item счёт пополняется на \term{Бонус запасы} --- инвесторы ликуют;
	\item R\&D возвращает откушенные мощности
	(\term{Нагрузка} снова прибавляется к мощности);
	\item и сверху --- \term{Бонус скорость}: кластер разгоняется навсегда.
\end{itemize}
Моментом внедрения считается \emph{время отправки посылки}, а не время
её проверки. Если тестирующая система задумалась --- это её галлюцинации,
а не ваши.

\item \textbf{Посылка на флажке.} Решение, отправленное до дедлайна,
но проверенное после, --- засчитывается. Отправленное после дедлайна ---
уходит в архив Комиссии по Выравниванию с пометкой
\emph{«надо было раньше»}.

\item \textbf{Без самодеятельности.} Сдавать некупленную задачу
бесполезно: прорыв, совершённый без бюджета, автоматика не внедряет,
а отправляет ведущему на разбирательство. Ведущий, конечно, может
и зачесть --- но сперва очень выразительно на вас посмотрит.

\item \textbf{Победа.} В момент дедлайна побеждает корпорация
с наибольшим счётом \term{Токенов}. Мощность, амбиции и красивые
питч-деки в зачёт не идут --- только токены на счету.

\item \textbf{Ведущий --- окончательный авторитет.} Он видит все сделки
в журнале, может продлить игру, исправить ошибочную запись и вручную
зачесть решение, если informatics прилёг отдохнуть. Спорить с ведущим
можно, но статистика споров не на вашей стороне: 0 побед из 0 возможных
у предыдущих составов.

\end{enumerate}

\section*{Параметры уровней}

Цены и бонусы зависят от уровня сложности задачи и объявляются
на табло игры (таблица \emph{«Параметры игры»}). Чем сложнее прорыв,
тем дороже утечка весов и тем щедрее инвесторы.

\begin{center}
\begin{tabular}{lccc}
\toprule
 & \textbf{Уровень 1} & \textbf{Уровень 2} & \textbf{Уровень 3} \\
\midrule
Цена задачи        & 12\,000 & 12\,000 & 12\,000 \\
Цена подсказки     & 3\,000  & 7\,000  & 10\,000 \\
Нагрузка от задачи & 2       & 2       & 2       \\
Бонус запасы       & 12\,000 & 25\,000 & 50\,000 \\
Бонус скорость     & 4       & 7       & 11      \\
\bottomrule
\end{tabular}

\small\emph{(значения по умолчанию; точные параметры вашей игры --- на табло)}
\end{center}

\section*{Шпаргалка по цветам табло}

\begin{center}
\begin{tabular}{ll}
\toprule
\colorbox{gray!30}{\strut ~серая~}   & задача не куплена: постановка в сейфе, условие скрыто \\
\colorbox{cyan!40}{\strut ~голубая~} & куплена: R\&D работает, кластер страдает \\
\colorbox{green!40}{\strut ~зелёная~} & решена: инвесторы аплодируют стоя \\
\bottomrule
\end{tabular}
\end{center}

\vspace{1em}
\begin{center}
\emph{Удачной генерации токенов! И помните: возвратов нет,
это записано в лицензионном соглашении, которое никто не читал.}
\end{center}

\end{document}
